% А2.2 — Корни и степень с рациональным показателем
% Старая методичка: стр. 65 (старое задание 7)
% Источники: методичка «Степени» (определение 4, замечание про отрицательные основания,
% свойство 10, задача 17), шпаргалка «Степень». Стиль — стиль-изложения.md

\begin{tema}{А2.2}
\opr{Корнем $n$-й степени} ($n\geqslant2$ — натуральное) из числа $a$ называется число,
$n$-я степень которого равна $a$. Обозначение: $\sqrt[n]a$.
\begin{itemize}[itemsep=0pt, topsep=2pt]
  \item \emph{$n$ чётное} (в том числе квадратный корень, $n=2$, пишут $\sqrt a$):
        корень извлекают только из $a\geqslant0$, и $\sqrt[n]a$ — это
        \emph{неотрицательное} число. Например, $\sqrt9=3$, а не $\pm3$; $\sqrt[4]{16}=2$.
  \item \emph{$n$ нечётное} (в том числе кубический, $n=3$): корень есть у любого числа.
        Например, $\sqrt[3]{-8}=-2$.
\end{itemize}

\begin{pfTablica}{lll}
& $n$ чётное & $n$ нечётное\\
Когда верно & $a\geqslant0$, $b\geqslant0$ & $a$, $b$ — любые\\
Произведение & \Block{1-2}{$\sqrt[n]a\cdot\sqrt[n]b=\sqrt[n]{ab}$}\\
Частное & $\dfrac{\sqrt[n]a}{\sqrt[n]b}=\sqrt[n]{\dfrac ab}$, \ $b>0$ & то же, $b\ne0$\\
Корень из степени & $\sqrt[n]{a^n}=|a|$ \ (любое $a$) & $\sqrt[n]{a^n}=a$\\
Степень корня & \Block{1-2}{$\left(\sqrt[n]a\right)^n=a$ \ ($a\geqslant0$ при чётном $n$)}\\
\end{pfTablica}

\opr{Вынесение множителя из-под корня:} $\sqrt{a^2b}=|a|\sqrt b$, например
$\sqrt{50}=\sqrt{25\cdot2}=5\sqrt2$. Обратно — \opr{внесение под корень}:
$a\sqrt b=\sqrt{a^2b}$ при $a\geqslant0$. Так удобно сравнивать:
$3\sqrt2=\sqrt{18}<\sqrt{20}=2\sqrt5$.

\opr{Иррациональность в знаменателе} убирают, домножая на сопряжённое выражение:
$\dfrac1{\sqrt a-\sqrt b}=\dfrac{\sqrt a+\sqrt b}{a-b}$. Например,
$\dfrac1{\sqrt3-\sqrt2}=\dfrac{\sqrt3+\sqrt2}{3-2}=\sqrt3+\sqrt2$.

\opr{Степенью} числа $a>0$ с рациональным показателем $\frac mn$ ($m$ — целое,
$n\geqslant2$ — натуральное) называется число $a^{\frac mn}=\sqrt[n]{a^m}=\left(\sqrt[n]a\right)^m$. Например, $8^{\frac23}=\left(\sqrt[3]8\right)^2=4$. При $a>0$, $b>0$ свойства из
\tref{А2.1} верны для любых показателей. Поэтому выражения с корнями удобно
переписывать степенями.

Важно отметить, что основание дробной степени должно быть положительным. Иначе
получилось бы $(-8)^{\frac13}=\sqrt[3]{-8}=-2$, но $\frac13=\frac26$ и
$(-8)^{\frac26}=\sqrt[6]{64}=2$. Поэтому пишут $\sqrt[3]{-8}$, а не $(-8)^{\frac13}$.

\begin{primer}
Найдите значение выражения $3\cdot\sqrt[4]{27}\cdot\sqrt[12]{27}$.
Решение: запишем корни как степени тройки: $\sqrt[4]{27}=3^{\frac34}$,
$\sqrt[12]{27}=3^{\frac3{12}}=3^{\frac14}$. Тогда
$3^{1+\frac34+\frac14}=3^2=9$. Ответ: 9.
\end{primer}

\begin{oshibka}
Корень из суммы не равен сумме корней: $\sqrt{9+16}=5$, а $\sqrt9+\sqrt{16}=7$.
\end{oshibka}
\end{tema}
